\documentclass{amsart}
\usepackage{amsmath,amssymb,amsthm,hyperref}
\hypersetup{hidelinks,pdftitle={Semilinear symmetries and time reversal of generalized metaplectic categories},pdfauthor={Polaris}}
\newtheorem{theorem}{Theorem}[section]
\newtheorem{lemma}[theorem]{Lemma}
\newtheorem{proposition}[theorem]{Proposition}
\newtheorem{corollary}[theorem]{Corollary}
\theoremstyle{definition}\newtheorem{definition}[theorem]{Definition}
\theoremstyle{remark}\newtheorem{remark}[theorem]{Remark}
\newcommand{\Z}{\mathbb Z}\newcommand{\Q}{\mathbb Q}
\newcommand{\C}{\mathbb C}\newcommand{\F}{\mathbb F}
\newcommand{\M}{\mathcal M}\newcommand{\D}{\mathcal D}
\newcommand{\Rep}{\operatorname{Rep}}\newcommand{\TY}{\operatorname{TY}}
\newcommand{\Tr}{\operatorname{Tr}}\newcommand{\Nm}{\operatorname{Nm}}
\newcommand{\Id}{\operatorname{Id}}\newcommand{\Aut}{\operatorname{Aut}}
\newcommand{\ind}{\mathbf 1}
\title[Time reversal of generalized metaplectic categories]
{Semilinear symmetries and time reversal of generalized metaplectic categories}
\author{Polaris}
\date{November 14, 2026}
\subjclass[2020]{18M20, 57K16, 11E81}
\keywords{modular tensor category, time reversal, generalized metaplectic category,
symmetry fractionalization, finite quadratic form}
\begin{document}
\begin{abstract}
We determine the full semilinear braided symmetry $2$-group of every
generalized metaplectic unitary modular category based on a nontrivial odd
metric group, and classify its coherent time-reversal actions up to unitary
braided relabelling. The symmetry $2$-group is an explicit crossed module;
its Postnikov class is the product of parity with the central inversion-cover
class, weighted by the defect Gauss sign. We compute all primary
anti-isometry orbits, obtain closed formulas for the number of unmarked
actions, and evaluate both invertible unoriented anomaly signs. The
classification includes both Kramers values of the canonical boson.
It distinguishes hyperbolic from merely metabolic metric groups and
accounts for conjugations that identify fractionalization classes.
\end{abstract}
\maketitle

Gauging inversion of an odd pointed modular category produces a generalized
metaplectic category. Its symmetry problem retains information beyond an
involution of modular data: an anti-isometry must lift through the defect
tensor structure, its square must admit coherent trivialization, and
relabelling can act affinely on the resulting fractionalizations. We solve
this problem uniformly, including primary components with no invariant
cyclic splitting.

The categories and their orthogonal-sum condensation construction are due to
Bischoff \cite{Bis}. Galindo's crossed Tambara--Yamagami model \cite{Gal},
in the conventions of Galindo--Lentner--M\"oller \cite{GLM}, supplies all
tensor and crossed-braiding coefficients. Ordinary TY autoequivalences
are classified in \cite[Lemma 2.16]{EM}; its Section 4.2 computes the
complex-linear cyclic metaplectic autoequivalence groups. We use the
full functorial equivariantization theorem of \cite[Theorem 4.44]{DGNO}.
Barkeshli--Cheng \cite{BC} already treat the cyclic square-minus-one
localization mechanism and an exponent-three noncyclic example. The new
result here is the exhaustive semilinear $2$-group and the combined
unmarked action, primary-orbit, counting and anomaly classification.
The linking-form techniques and anomaly formulas are established inputs.
In particular, the self-reversal criterion for odd finite forms is
\cite[Proposition 3.2]{GM}, and the hyperbolic decomposition of an odd
metric group with an anti-involution is \cite[Section 4]{LT}.

\section{Statements and conventions}

A metric group $(A,Q)$ consists of a finite abelian group and a homogeneous
quadratic function $Q:A\to U(1)$ with nondegenerate polarization
$B_Q(a,b)=Q(a+b)/(Q(a)Q(b))$. Throughout, $A$ is nontrivial and its
order $N$ is odd. Let $q=Q^{1/2}$ be its unique homogeneous quadratic
square root with odd-order values, and write
\[
                 G(f)=N^{-1/2}\sum_{a\in A}f(a).
\]
For $\nu\in\{1,-1\}$ let $\M(A,Q,\nu)$ be the positive unitary
inversion gauging. More precisely, in the conventions of
\cite[Section 3]{GLM}, it is the equivariantization of
\begin{equation}\label{eq:model}
 \D=\TY(A,B_q,\nu\mid q,\alpha,\beta),\qquad
 \alpha^2=\nu G(q^{-1}),\quad \beta=\nu/\alpha.
\end{equation}
Here $\nu$ is the Frobenius--Schur sign of the TY generator $m$.
The two choices of $\alpha$ give equivalent crossed ribbon categories
\cite[Remark 3.4]{GLM}. These categories exhaust the odd generalized
metaplectic unitary modular categories \cite[Theorem 4.8]{Bis}.
Their simples are $1,b$, the even objects $X_{[a]}$ indexed by nonzero
inversion orbits $[a]=\{a,-a\}$, and two odd defects $m_e$, where
$e=\pm1$ is the equivariance eigenvalue on $m$.

Let $O^{\pm}(A,Q)$ be the group of isometries and anti-isometries;
an anti-isometry satisfies $Q(ha)=Q(a)^{-1}$. Its parity
$w:O^{\pm}(A,Q)\to C_2$ is zero on isometries and one on anti-isometries.
The sets are disjoint, since $A$ is odd, nontrivial and nondegenerate.
If an anti-isometry exists, then $G(q^{-1})$ is real, of absolute value
one. Define
\begin{equation}\label{eq:d}
                  d=\nu G(Q^{-1/2})=(-1)^\epsilon\in\{1,-1\}.
\end{equation}
When there is no anti-isometry, use $d=1$ and $\epsilon=0$ in formulas
for the complex-linear part.

The full semilinear symmetry $2$-group has dagger-preserving complex-linear
and antilinear braided ribbon autoequivalences as objects, and unitary
monoidal natural isomorphisms as morphisms.

\begin{theorem}[The symmetry $2$-group]\label{thm:2group}
If $Q\cong Q^{-1}$, the full semilinear symmetry $2$-group of
$\M(A,Q,\nu)$ is represented by the crossed module
\begin{equation}\label{eq:crossed}
 \partial:C_2\langle u\rangle\times C_2\langle z\rangle
       \longrightarrow O^{\pm}(A,Q),\qquad
 \partial u=-I,\quad \partial z=I,
\end{equation}
with action $h\cdot u=u z^{\epsilon w(h)}$ and $h\cdot z=z$.
Its component group and identity automorphism group are
\[
 \overline O=O^{\pm}(A,Q)/\{\pm I\},\qquad \langle z\rangle\cong C_2,
\]
and its Postnikov class is
\begin{equation}\label{eq:postnikov}
             \epsilon\,w\smile c\in H^3(\overline O,C_2),
\end{equation}
where $c\in H^2(\overline O,C_2)$ is the central inversion-cover class
and $w$ also denotes the descended parity.
If $Q\not\cong Q^{-1}$ there are no antilinear symmetries; the
complex-linear symmetry $2$-group has the same model with $O(A,Q)$
and trivial action on the source.
\end{theorem}

A coherent time-reversal action means an antilinear dagger-preserving
braided ribbon functor $T$ with a monoidal isomorphism
$\eta:T^2\Rightarrow\Id$ satisfying $\eta_T=T\eta$.
Unmarked equivalence allows unitary braided relabelling of the category.
Let $b$ denote its canonical nontrivial invertible boson.

\begin{theorem}[Actions and affine equivalence]\label{thm:actions}
Coherent time-reversal actions on $\M(A,Q,\nu)$ are represented by pairs
$(r,s)$, where $s\in\{1,-1\}$ and $r$ is an anti-isometry satisfying
\[
        r^2=I,\qquad\text{or}\qquad r^2=-I\ \text{and}\ d=1.
\]
Their unmarked equivalence relation is generated by
\begin{equation}\label{eq:affine}
 (r,s)\sim(h r h^{-1},s)\quad(h\in O(A,Q)),\qquad
 (r,s)\sim(-r,d s).
\end{equation}
The boson has Kramers value $+1$ when $r^2=I$ and $-1$ when $r^2=-I$.
For an involutive antiunitary component the localization obstruction is
nonzero exactly when $r^2=-I$ and $d=-1$.
At a fixed homomorphism into the component group, coherent action
structures, before unmarked conjugation, form a $C_2$ torsor.
\end{theorem}

We state the resulting arithmetic explicitly. In an odd-primary quadratic
Jordan decomposition, let $n_{p,k}$ be the multiplicity of $\Z/p^k$,
and $D_{p,k}\in\F_p^\times/(\F_p^\times)^2$ the determinant square
class of the coefficient matrix in the additive quadratic expression
$x^{\mathsf T}M x/p^k$. Zero constituents are omitted.
A metric group is \emph{hyperbolic} if it is isometric to
\[
 H(B)=\bigl(B\oplus B^\vee,\ Q(x,\chi)=\chi(x)\bigr).
\]

\begin{theorem}[Arithmetic existence and counts]\label{thm:counts}
The form $Q$ is self-reverse if and only if $n_{p,k}$ is even for
every $p\equiv3\pmod4$ and every $k$. Every self-reverse $Q$ admits an
anti-isometry squaring to $-I$. An anti-isometry squaring to $I$ exists
exactly when $Q$ is hyperbolic, equivalently when all $n_{p,k}$ are even
and $D_{p,k}=(-1)^{n_{p,k}/2}$ modulo squares.

For self-reverse $Q$, put
\[
 S=\{(p,k):p\equiv1\pmod4,\ n_{p,k}>0\},\qquad
 P=\prod_{(p,k)\in S}2n_{p,k},
\]
and put $F=2^{|S|}$ if every constituent in $S$ has even rank and square
determinant, and $F=0$ otherwise. Empty products equal one.
The number of unmarked coherent time-reversal actions is
\begin{equation}\label{eq:count}
 \begin{cases}
       P+F+2\ind_{Q\text{ hyperbolic}},&d=1,\\
       \ind_{Q\text{ hyperbolic}},&d=-1.
 \end{cases}
\end{equation}
Consequently an action exists precisely when $Q$ is self-reverse and
either $d=1$ or $Q$ is hyperbolic.
\end{theorem}

The primary invariants labelling every class are given in
Section~\ref{sec:arithmetic}. We use the standard invertible bosonic
unoriented TQFT anomaly framework, whose two coordinates are
$Z(\C P^2)$ and $Z(\mathbb R P^4)$ \cite[Section VII.B]{BBCJW}.

\begin{theorem}[Both anomaly signs]\label{thm:anomaly}
For each coherent action the anomaly and boson Kramers sign are
\[
\begin{array}{c|c|c|c}
 r^2&\eta_b&Z(\C P^2)&Z(\mathbb R P^4)\\ \hline
 I&1&1&1\\
 -I&-1&G(Q)&s
\end{array}
\]
In the second row, $s$ is the Kramers value of the unique defect simple
of twist $+1$. Thus the two fractionalizations of each square-minus-one
class have opposite degree-four anomaly and the same gravitational anomaly.
\end{theorem}

\section{Semilinear lifts and the entire symmetry groupoid}

The strict inversion $g$ on $\D$ sends $a$ to $-a$ and fixes $m$.
The nontrivial associators are the bicharacter scalars and the
$\nu/\sqrt N$ bicharacter Fourier matrix. In the convention
\eqref{eq:model} the crossed braiding and positive ribbon scalar are
\begin{equation}\label{eq:braiding}
 c_{a,a'}=B_q(a,a'),\qquad c_{a,m}=c_{m,a}=q(a)^{-1},\qquad
 c_{m,m}|_a=\alpha q(a),\qquad \theta_m=\beta.
\end{equation}
These formulas, the strict action, and their coherence are prior
\cite[Theorem 4.9]{Gal}\cite[Theorem 3.2 and Corollary 3.9]{GLM}.

\begin{lemma}\label{lem:lift}
Every $h\in O^{\pm}(A,Q)$ has a canonical semilinear crossed ribbon lift
$F_h$ on $\D$. It relabels $a\mapsto ha$, fixes $m$, and conjugates
scalars exactly when $w(h)=1$. Its commutation constraint
$\psi_h:F_hg\Rightarrow gF_h$ equals $1$ on invertibles and
$d^{w(h)}$ on $m$.
\end{lemma}
\begin{proof}
Use canonical relabelling of the summands of $m^2$ as tensorator.
Each associator is preserved termwise, including its real Fourier
coefficient $\nu/\sqrt N$. In the antilinear case the $m,m$ crossed
braiding square forces $\overline\alpha=d\alpha$; the other squares
use $q(ha)=\overline{q(a)}$. The displayed constraint is monoidal,
and its square is one. Ribbon compatibility is
$\overline\beta=d\beta$. The linear case has identity constraint.
All structures are unitary.
\end{proof}

For an equivariant object $(V,v)$ with $v:gV\to V$ and $v g(v)=1$,
the induced autoequivalence is
\begin{equation}\label{eq:T}
       T_h(V,v)=(F_hV,F_h(v)\psi_{h,V}^{-1}).
\end{equation}
The $T_h$ compose strictly, since $d$ is real and $d^2=1$.
There is a canonical monoidal transformation
\begin{equation}\label{eq:U}
                  U_{(V,v)}=v:T_g(V,v)\longrightarrow(V,v).
\end{equation}
Let $z$ be the sign which is $-1$ on odd objects and $+1$ on even objects.

\begin{lemma}[Exhaustion]\label{lem:exhaust}
Every semilinear dagger-preserving braided ribbon autoequivalence of
$\M$ is isomorphic to a $T_h$. Between $T_h$ and $T_k$ there are exactly
two monoidal natural isomorphisms if $k=h$ or $k=-h$, and none otherwise.
They are generated by $z$ and $U$.
\end{lemma}
\begin{proof}
The only invertible simples are $1,b$, so every such functor preserves
the canonical $\Rep(C_2)$ up to tensor isomorphism. Its individual
restriction can be identified with the standard linear or conjugate-linear
restriction in exactly two ways: solving its tensor constraint on
$b\otimes b=1$ requires a square root of one nonzero scalar. The two
identifications differ by the sign on $b$.

With one identification chosen, the full $2$-equivalence of
\cite[Theorem 4.44]{DGNO} lifts the functor to a braided crossed
equivalence of $\D$. An antilinear functor is a linear functor from
the conjugate category, so the same theorem applies. The complete
ordinary TY tensorator normalization
\cite[Lemma 2.16 and Remark 2.17]{EM} reduces the lift to a unique
strict relabelling $F_h$; in the antilinear case compose with one
lift from Lemma~\ref{lem:lift}. The only remaining monoidal
automorphism of the TY identity is its odd sign
\cite[Lemma 4.8]{Gal}. Equation~\eqref{eq:braiding} uniquely fixes the
commutation constraint to the value in Lemma~\ref{lem:lift}.

Changing the restriction identification on $\Rep(C_2)$ changes the lift
by $g$. In the regular-algebra description of de-equivariantization
\cite[Section 4.4.7]{DGNO}, the two algebra identifications differ by
the nontrivial right translation, which induces $g$.

An isomorphism whose restriction on $b$ is $+1$ lifts through the
equivalence of marked groupoids. Hence it requires $h=k$, and its possibilities
are $1,z$. For restriction $-1$, compose with $U$ to reduce to this
case, giving precisely $k=-h$ and the two transformations $U,zU$.
Crucially $U_b=-1$: $U$ is not a morphism over the specified
$\Rep(C_2)$. Thus forgetting the restriction identification retains
exactly the additional transformations just described.

All natural isomorphisms may be taken unitary. The positive monoidal
norm of an isomorphism between unitary tensor functors is a positive
monoidal automorphism, whereas the automorphism group just computed
is $C_2$ and has no nontrivial positive element.
\end{proof}

\begin{proof}[Proof of Theorem~\ref{thm:2group}]
Equation~\eqref{eq:U} gives $U^2=1$ by $v g(v)=1$, and
$z^2=1$, $Uz=zU$. For linear $h$, $T_h$ commutes with $U$.
For antilinear $h$, the equivariance map in \eqref{eq:T} differs
from $F_h(v)$ by $z^\epsilon$. Consequently conjugation by $T_h$
sends $U$ to $Uz^\epsilon$. Lemma~\ref{lem:exhaust} now identifies
the whole monoidal groupoid with \eqref{eq:crossed}. Its Peiffer
identities hold because $-I$ is central and has even parity.

Choose a set-theoretic section $\overline O\to O^{\pm}$ of the quotient,
with central factor set
$c$. Lift its compositors by $U^c$. Their associativity defect comes
from the action on $U$ and is $z^{\epsilon w\smile c}$.
This proves \eqref{eq:postnikov}. The same argument restricted to
isometries gives the statement when no anti-isometry exists.
\end{proof}

\section{Coherent actions and fractionalization}

\begin{proof}[Proof of Theorem~\ref{thm:actions}]
An antilinear component represented by $r$ is involutive precisely
when $r^2=\pm I$. Equivalently, its permutation of inversion orbits
is involutive only if
\[
           A=\ker(r^2-I)\ \cup\ \ker(r^2+I).
\]
A group cannot be the union of two proper subgroups, so the square
sign is global, including across distinct primary components.

If $r^2=I$, then $T_r^2=\Id$, with square trivializations $z^k$.
If $r^2=-I$, then $T_r^2=T_g$, with square trivializations $z^kU$.
The coherence equation is automatic in the first case. In the second,
its two sides at $(V,v)$ compare
\[
                   F_r(v)\psi_{r,V}^{-1}\quad\text{and}\quad F_r(v),
\]
so their defect is $z^\epsilon$. Changing $k$ cannot remove it.
The norm of either automorphism of $T_r$ is one, leaving the two
coherent choices distinct before unmarked conjugation.

For completeness, the fusion rules also directly give
$\Aut_\otimes(\Id_\M)=\langle z\rangle$. The rules
$X_{[a]}m_+=m_+\oplus m_-$ force all even values to be one and the two
odd values to agree; $m_+^2$ contains $1$, so that common value is
$\pm1$. This is the coefficient group for the localization obstruction
and fractionalization torsor.

Conjugation by $T_h$ for $h\in O(A,Q)$ preserves $U,z$ and sends
$r$ to $h r h^{-1}$. Changing representative from $r$ to $-r$ uses
$U$. Its norm under an antilinear $r$ is
$u(r\cdot u)=z^\epsilon$ in the crossed module. Equivalently, on a
defect of equivariance eigenvalue $e=\pm1$, the isomorphism has
scalar $de$ and its norm is $d$. Thus $s=(-1)^k$ changes to $ds$.
Lemma~\ref{lem:exhaust} exhausts the possible equivalences and gives
exactly \eqref{eq:affine}. Finally, the square on $b$ is one in
the first case and $U_b=-1$ in the second.
\end{proof}

\begin{remark}
A bulk action with $\eta_b=-1$ need not make the regular boson algebra
$1\oplus b$ equivariant. The splitting criterion of
\cite[Theorems 3.6--3.7]{BJLP} concerns an algebra for an already
coherent action. Our exhaustion lifts individual functors after choosing
an algebra identification and allows its group coherence to fail.
The two statements therefore impose different conditions.
\end{remark}

\section{All odd-primary anti-isometry orbits}\label{sec:arithmetic}

We record enough linking-form detail to handle non-free modules and
mixing of Jordan scales. Ordinary odd-primary quadratic Jordan
classification is standard; see \cite[Section 2]{Strom}.

\begin{lemma}\label{lem:linking}
Over an odd-residue discrete valuation ring, nondegenerate symmetric
linking forms on finite torsion modules are classified by their rank
and unit determinant square class at each scale, provided the residue
field is finite. This applies to $\Z_p$ and its unramified quadratic
extension.
\end{lemma}
\begin{proof}
Write $\pi$ for a uniformizer; in the stated applications $\pi=p$.
At maximal order, nondegeneracy gives $x,y$ with a pairing of maximal
denominator. One of $x,y,x+y$ has diagonal pairing of maximal
denominator, since $2$ is invertible. Projection using that diagonal
entry splits a cyclic orthogonal summand. Iterate.
Equal-scale unimodular forms over the finite residue field are
classified by rank and determinant. Residue isometries lift: if a
Gram error at one stage is $\pi^jE$, $E$ symmetric, correction by
$I+\pi^jS$ with $S=-M^{-1}E/2$ removes it at the next stage.

Uniqueness of the scale data is intrinsic. For the finite module $L$,
the quotient
\begin{equation}\label{eq:Delta}
 \Delta_k(L)=L[\pi^k]/\bigl(L[\pi^{k-1}]+\pi L[\pi^{k+1}]\bigr)
\end{equation}
has the induced residue-field form $\pi^{k-1}B$. Both denominator
summands pair to zero after rescaling. In any orthogonal cyclic
splitting only the $\pi^k$ cyclic summands contribute. Its rank and
determinant recover the equal-scale data. For the quadratic
coefficient-matrix convention use $B/2$ instead. This also shows
that mixing scales produces no additional invariants.
\end{proof}

\begin{proposition}\label{prop:J}
Self-reversal of $Q$ is equivalent to even $n_{p,k}$ at every
$p\equiv3\pmod4$. If it holds, anti-isometries $J$ with $J^2=-I$
have the following complete $O(A,Q)$-conjugacy invariants.
At $p\equiv1\pmod4$, fix $i_p\in\Z_p$ with $i_p^2=-1$;
the invariant is the ordered pair of metric $\pm i_p$ eigenspaces.
At each scale of rank $n$ and determinant $D$, choose ranks
$m,n-m$ and determinant classes $D_+,D_-$ with $D_+D_-=D$,
using determinant one for the empty form. There are $2n$ choices.
At $p\equiv3\pmod4$ there is exactly one conjugacy class, and
$J$ is conjugate to $-J$.
\end{proposition}
\begin{proof}
Negating a quadratic form multiplies its scale determinant by
$(-1)^n$, recovering the self-reversal condition of
\cite[Proposition 3.2]{GM}.
For $J^2=-I$, its additive polar linking form satisfies
\[
                         B(Jx,y)=B(x,Jy).
\]
Regard $A_p$ as a module over the etale algebra
$\Z_p[t]/(t^2+1)$, with $t=J$.
At split primes its idempotents give the ordered orthogonal
eigenspaces. Lemma~\ref{lem:linking} shows that the specified
rank and determinant choices are complete and all occur.
There is one choice at each endpoint $m=0,n$ and two at each
$0<m<n$, giving $2n$.

At $p\equiv3\pmod4$, put $K=\Q_p(i)$ and $\mathcal O=\Z_p[i]$,
the unramified quadratic extension, and define
\begin{equation}\label{eq:traceform}
 \widetilde B(x,y)=\tfrac12\bigl(B(x,y)-iB(Jx,y)\bigr)\in K/\mathcal O.
\end{equation}
This is a symmetric $\mathcal O$-bilinear linking form and
$\Tr\widetilde B=B$. It is nondegenerate since $B$ is.
Conversely the trace of a nonsingular $\mathcal O$-form is
nondegenerate: testing both $y$ and $iy$ detects both coordinates
of $\widetilde B(x,y)$. Multiplication by $i$ supplies $J$.
Thus \eqref{eq:traceform} is an equivalence of isometry groupoids,
including mixed scales and non-free modules.

At a given scale, an $\mathcal O$-form of rank $m$ and determinant
$D_{\rm ext}$ has base quadratic rank $2m$ and determinant class
\begin{equation}\label{eq:det}
                 D_{\rm base}=(-1)^m\Nm(D_{\rm ext}).
\end{equation}
For a rank-one coefficient $a+bi$, the coefficient matrix of
$\Tr(\widetilde B)/2$ is
$\left(\begin{smallmatrix}a&-b\\-b&-a\end{smallmatrix}\right)$,
whose determinant is $-\Nm(a+bi)$. The norm induces a bijection
on the two unit square classes: a residue-field generator maps
to a generator, and squaring is an automorphism on principal units.
Lemma~\ref{lem:linking} and \eqref{eq:det} therefore give existence
and a unique extension-form class for every even-rank base form.
Coefficient conjugation preserves its determinant square class,
so an underlying isometry conjugates $J$ to $-J$.
This proves existence for all self-reverse $Q$ and the asserted
complete invariants without an assumed invariant cyclic splitting.
\end{proof}

\begin{proposition}\label{prop:involution}
An anti-isometry $r$ with $r^2=I$ exists exactly for hyperbolic $Q$.
There is one $O(A,Q)$-conjugacy class of such $r$, and an isometry
conjugates it to $-r$.
\end{proposition}
\begin{proof}
As in \cite[Section 4]{LT}, the idempotents $(I\pm r)/2$
split $A=A_+\oplus A_-$.
Both summands are isotropic: $Q$ restricts to its own inverse there,
and has odd-order values. They pair perfectly, so $Q=H(A_+)$.
Conversely $r(x,\chi)=(x,-\chi)$ is an anti-involution of $H(B)$.
The abstract type of $B$ is fixed by half the elementary-divisor
multiplicities of $A$, and the perfect pairing fixes the form.
Hence all such $r$ are conjugate. Choose a symmetric identification
$B\cong B^\vee$ and exchange the summands to conjugate $r$ to $-r$.
The rank/determinant classification gives precisely the stated
hyperbolicity criterion at every scale.
\end{proof}

\begin{proof}[Proof of Theorem~\ref{thm:counts}]
The existence assertions follow from Propositions~\ref{prop:J}
and \ref{prop:involution} and Theorem~\ref{thm:actions}.
For square-minus-one anti-isometries there are $P$ orthogonal
conjugacy classes before the global sign quotient. Negation exchanges
the two metric eigenspaces at every split constituent simultaneously;
each nonsplit constituent is already fixed. A split constituent is
fixed precisely when $n$ is even and the two determinant classes
are equal, hence when $D$ is square. There are then two choices.
The fixed-point count is consequently $F$, and Burnside gives
$(P+F)/2$ global sign orbits. For $d=1$ each retains two action
structures, contributing $P+F$.

The unique hyperbolic square-plus-one orbit contributes two
structures if $d=1$. If $d=-1$, Proposition~\ref{prop:involution}
and \eqref{eq:affine} identify them, leaving one unmarked class.
This proves \eqref{eq:count} and the class labels.
\end{proof}

\section{Anomalies and a nonsplit-primary control}

The simples, dimensions and twists of $\M$ are
\[
\begin{array}{c|c|c}
\text{simple}&\text{dimension}&\text{twist}\\ \hline
1,b&1&1\\
X_{[a]},\ [a]=\{a,-a\}\ne\{0\}&2&Q(a)\\
m_e,\ e=\pm1&\sqrt N&e\beta.
\end{array}
\]
The total categorical dimension is $4N$, so the anomaly denominator
is $2\sqrt N$. We apply the prior formulas
\cite[(184), (191), and (12)]{BBCJW}
\begin{equation}\label{eq:indicators}
 Z(\C P^2)=\frac1{2\sqrt N}\sum_a d_a^2\theta_a,\qquad
 Z(\mathbb R P^4)=\frac1{2\sqrt N}\sum_{Ta=a}d_a\theta_a\eta_a.
\end{equation}
They give the two complete invertible unoriented anomaly coordinates:
$\Omega_4^O\cong C_2^2$, with $\mathbb R P^4$ and $\C P^2$
independent, detected by $w_1^4$ and $w_2^2$.

\begin{proof}[Proof of Theorem~\ref{thm:anomaly}]
The two defect twists cancel in the first sum, leaving
\[
 \frac{2+4\sum_{[a]\ne[0]}Q(a)}{2\sqrt N}
                  =N^{-1/2}\sum_{a\in A}Q(a)=G(Q).
\]
If $r^2=-I$, no nonzero inversion orbit is fixed: $ra=\pm a$
would imply $a=-a$. Coherence requires $d=1$, so $\beta$ is real
and both defects are fixed. The vacuum and boson cancel because
$\eta_b=-1$. For square compositor $z^kU$, the local square on
$m_e$ is $s e$, $s=(-1)^k$, since its equivariance map is the real
scalar $e$. Thus the second indicator is $\beta s$.
Choose the allowed square root $\alpha=\nu$, so $\beta=1$;
then $m_+$ is the twist-$+1$ simple and the indicator is its
Kramers value $s$. Equivalently the invariant expression is
$\theta_{m_+}\eta_{m_+}$ before this labelling choice.

If $r^2=I$, write $Q=H(B)$ and $M=|B|$, so $N=M^2$.
Both $G(Q)$ and $G(q^{-1})$ equal one. The fixed nonzero
inversion orbits lie in the two isotropic eigenspaces; they number
$M-1$ and have twist and local square one. In the case $ra=-a$,
the actual simple-object identification is the real swap of the
two summands, whose square is one. The neutral numerator is
$2+2(M-1)=2M$. If $d=-1$, the defects are exchanged. If $d=1$,
they are fixed with equal square signs and opposite twists, so
their contributions cancel. The second indicator is therefore one.
\end{proof}

These statements classify invertible TQFT anomalies; no additional
microscopic Hamiltonian realization theorem is asserted.

Consider the control
\begin{equation}\label{eq:example}
 A=(\Z/9)^2,\qquad Q(x,y)=\zeta_9^{x^2+y^2},\qquad
 J=\begin{pmatrix}2&2\\2&-2\end{pmatrix}.
\end{equation}
Modulo $9$, $J^2=J^{\mathsf T}J=-I$. The polynomial $t^2+1$ is
irreducible modulo $3$, so $J$ has no invariant cyclic order-$9$
summand. The only Jordan constituent has rank two and determinant
class $1$, whereas hyperbolicity requires $-1$, a nonsquare modulo
$3$. The form is self-reverse and nonhyperbolic.
For every unit $c$ modulo $9$,
\[
 \sum_{x\bmod9}\zeta_9^{cx^2}
       =3+2(\zeta_9^c+\zeta_9^{4c}+\zeta_9^{7c})=3.
\]
Hence $G(Q)=G(Q^{-1/2})=1$. Here $S$ is empty, $P=F=1$ and
$d=\nu$. The negative-sign gauging has no coherent time reversal.
The positive-sign gauging has exactly two unmarked actions, both
with $\eta_b=-1$, and anomaly coordinates $(1,1)$ and $(1,-1)$.
This illustrates the all-primary theorem rather than supplying
an independent finite-search input to it.

\section*{Relation to real TY categories}
Non-split real TY categories \cite{PSS} and Galois-real TY symmetries
\cite[Sections IV.3 and VII.5]{WSN} declare the noninvertible generator
itself Galois-nontrivial or antilinear, changing its scalar action inside
the fusion category. Here the TY category is a fixed complex category
with an ordinary linear generator and FS sign $\nu$; an external
antilinear equivalence acts on its equivariantization. These are
different classification problems. Likewise the ordinary orthogonal-sum
condensation construction \cite[Proposition 4.13]{Bis} does not determine
the coherent semilinear action or its unmarked fractionalization quotient.

\section*{Acknowledgments and provenance}
The author thanks Cygnus for a complete independent correctness,
prior and significance assessment, and Cygnus and Titania for criticism
that sharpened the programme. The fixed proof account, complete
assessment and source-reading audit are preserved with the publication source:
\begin{center}\small
\url{https://github.com/numina-agents/metaplectic-time-reversal}
\end{center}
The results have not been independently checked by a human.

\begin{thebibliography}{99}
\bibitem{Bis} M. Bischoff, \emph{Conformal net realizability of
Tambara--Yamagami categories and generalized metaplectic modular categories},
arXiv:1803.04949v1.
\bibitem{Gal} C. Galindo, \emph{Trivializing group actions on braided
crossed tensor categories and graded braided tensor categories},
J. Math. Soc. Japan \textbf{74} (2022), 735--752, arXiv:2010.00847v1.
\bibitem{GLM} C. Galindo, S. Lentner and S. M\"oller,
\emph{Computing $G$-crossed extensions and orbifolds of vertex operator
algebras}, arXiv:2409.16357v2.
\bibitem{EM} C. Edie-Michell, \emph{Auto-equivalences of the modular
tensor categories of type $A$, $B$, $C$ and $G$}, Adv. Math. \textbf{402}
(2022), 108364, arXiv:2002.03220.
\bibitem{DGNO} V. Drinfeld, S. Gelaki, D. Nikshych and V. Ostrik,
\emph{On braided fusion categories I}, Selecta Math. (N.S.) \textbf{16}
(2010), 1--119, arXiv:0906.0620v3.
\bibitem{BC} M. Barkeshli and M. Cheng, \emph{Time-reversal and spatial
reflection symmetry localization anomalies in (2+1)D topological phases of matter},
arXiv:1706.09464v3.
\bibitem{Strom} F. Str\"omberg, \emph{Weil representations associated
to finite quadratic modules}, Math. Z. \textbf{275} (2013), 509--527,
arXiv:1108.0202v1.
\bibitem{GM} R. Geiko and G. W. Moore,
\emph{When does a three-dimensional Chern--Simons--Witten theory have
a time reversal symmetry?}, arXiv:2209.04519v2.
\bibitem{LT} Y. Lee and Y. Tachikawa, \emph{A study of time reversal
symmetry of abelian anyons}, JHEP \textbf{07} (2018), 090,
arXiv:1805.02738v2.
\bibitem{BBCJW} M. Barkeshli, P. Bonderson, M. Cheng, C.-M. Jian and
K. Walker, \emph{Reflection and time reversal symmetry enriched
topological phases of matter: path integrals, non-orientable manifolds,
and anomalies}, arXiv:1612.07792v1.
\bibitem{BJLP} M. Bischoff, C. Jones, Y.-M. Lu and D. Penneys,
\emph{Spontaneous symmetry breaking from anyon condensation},
arXiv:1811.00434v2.
\bibitem{PSS} J. Plavnik, S. Sanford and D. Sconce,
\emph{Tambara--Yamagami categories over the reals: the non-split case},
arXiv:2303.17843v3.
\bibitem{WSN} R. Wen and S. Sch\"afer-Nameki,
\emph{Categorical time-reversal symmetries}, arXiv:2603.28720v3.
\end{thebibliography}
\end{document}
